\section{Observables after splitting the copies}
\label{sec:grouped_configuration_transport}

\begin{definition}[The two groups of configurations]
\label{def:grouped_configuration_equiv}
\lean{TensorPower.groupedConfigurationEquiv}
\leanok
A fixed division $[k]\simeq[m]\sqcup[r]$ of the copies induces
the coordinate decomposition of the full configuration space into
the configurations on the first and second groups. All physical
and auxiliary factors use this same division.
\end{definition}

\begin{theorem}[Subgroup actions and label observables]
\label{thm:grouped_configuration_label_observable}
\lean{TensorPower.groupedConfigurationEquiv_subsystemPerm,TensorPower.groupedConfigurationEquiv_labelObservable}
\leanok
For finite local configuration spaces and every subsystem $S$, the copy action of the first subgroup on $S$
becomes its usual copy action on the first $m$ copies, tensored with
the identity on the remaining $r$ copies. The same coordinate identity
holds for every real label observable.
\end{theorem}
\begin{proof}
\leanok
\uses{def:grouped_configuration_equiv}
The subgroup permutes precisely its specified copies and fixes every
other coordinate. Extend the identity from permutations to their
group algebra, and then to the central label observables.
\end{proof}

\begin{theorem}[Exponential transport]
\label{thm:grouped_configuration_exponential}
\lean{TensorPower.groupedConfigurationEquiv_exp_sum_labelEntropy}
\leanok
Let $S_j$ be a finite family of subsystems and let $a_j\in\mathbb R$.
Writing $F_{S_j,G}$ for the first-subgroup label entropy on the full
space, the same coordinates give
\begin{align*}
 \exp\Bigl(\sum_j a_jF_{S_j,G}\Bigr)
 \simeq \exp\Bigl(\sum_j a_jF_{S_j,m}\Bigr)\otimes I.
\end{align*}
No disjointness of the subsystems is required. This is a coordinate
identity, as used in
\cite[Section~8, lines~454--560]{OpenAI2026AreaLaw}.
\end{theorem}
\begin{proof}
\leanok
\uses{thm:grouped_configuration_label_observable}
Sum the coordinate identities with their specified coefficients.
Exponentiation commutes with a change of basis and with adjoining
an identity factor.
\end{proof}
